\documentclass[11pt]{article}
\usepackage{mathblatt}
% gesetzt von werkzeuge/setzer.py (Sorte selbst) aus dem Lernweg KUV
% und den Bankzeilen; Muster bau/proben/2026-10-09/hypotenuse-muster.
\usepackage{enumitem}
\newgeometry{a4paper,margin=18mm,top=15mm,bottom=20mm,footskip=9mm}
\pagestyle{fancy}\fancyhf{}\renewcommand{\headrulewidth}{0pt}
\fancyfoot[L]{\footnotesize\color{mbgrau}KUV \textperiodcentered{} Lösungen \textperiodcentered{} Wahrscheinlichkeit: günstig durch möglich}
\fancyfoot[R]{\footnotesize\color{mbgrau}\thepage}
\setlength{\parskip}{3pt}
\renewcommand{\kreuz}[1]{\mbox{$\square$\ #1}\quad}
\newcommand{\kreuzl}[1]{\par\hangindent1.4em\hangafter1\noindent$\square$\ #1}
\newcommand{\aufg}[3]{\par\noindent\makebox[0pt][r]{#3}\begin{minipage}[t]{\linewidth}#1\end{minipage}\par}
\newcommand{\aufgg}[4]{\par\noindent\makebox[0pt][r]{#4}\begin{minipage}[t]{0.6\linewidth}\vspace{0pt}#1\end{minipage}\hfill
  \begin{minipage}[t]{0.37\linewidth}\vspace{0pt}\centering #2\end{minipage}\par}
\newcommand{\nr}[1]{\makebox[7mm][l]{\textbf{#1}}}
\newcommand{\lsg}[3]{\par\noindent\begin{minipage}[t]{9mm}\textbf{#1}\end{minipage}%
  \begin{minipage}[t]{52mm}\raggedright\bfseries\boldmath #2\end{minipage}\hspace{3mm}%
  \begin{minipage}[t]{\dimexpr\linewidth-64mm\relax}\raggedright\small #3\end{minipage}\par\vspace{4pt}}
\newlength{\kr}
\begin{document}

\noindent{\Large\bfseries Lösungen}\hspace{1em}{\color{mbgrau}Wahrscheinlichkeit: günstig durch möglich}\par\smallskip
\noindent{\small Links das Ergebnis, rechts der Weg in kurzen Schritten. Stimmt dein Ergebnis nicht, such den ersten Schritt, der anders ist.}\par
\par\Needspace{25mm}\vspace{6pt}\noindent\colorbox{black!12}{\makebox[7mm]{\bfseries\strut A}}\hspace{2mm}{\bfseries Günstig durch möglich}\par\vspace{4pt}
\lsg{1.}{$P = \frac{1}{6}$}{möglich 6 (1 bis 6), günstig 1: $P = \frac{1}{6}$}
\lsg{2.}{$P = \frac{7}{12}$}{möglich $4 + 7 + 1 = 12$ Kugeln, günstig 7: $P = \frac{7}{12}$}
\lsg{3.}{$P = \frac{1}{2}$}{möglich 6 Felder, günstig 3 Felder mit A: $P = \frac{3}{6} = \frac{1}{2}$. Es zählen die Felder, nicht die Buchstaben.}
\lsg{4.}{$P = \frac{1}{5}$}{möglich 20, günstig 5, 10, 15, 20 (4 Zettel): $P = \frac{4}{20} = \frac{1}{5}$}
\lsg{5.}{Aus dem ersten: $\frac{3}{8} > \frac{1}{3}$}{Aus dem ersten Beutel. Erster Beutel: $P = \frac{3}{8} = 0{,}375$; zweiter Beutel: $P = \frac{4}{12} = \frac{1}{3} \approx 0{,}333$. Im zweiten sind zwar mehr rote Kugeln, aber ihr Anteil ist kleiner.}
\par\Needspace{25mm}\vspace{6pt}\noindent\colorbox{black!12}{\makebox[7mm]{\bfseries\strut B}}\hspace{2mm}{\bfseries Alle mitzählen, in Prozent}\par\vspace{4pt}
\lsg{1.}{$0{,}6 = 60\,\%$}{$3 : 5 = 0{,}6 = 60\,\%$}
\lsg{2.}{$\frac{1}{10} = 10\,\%$}{alle $54 + 6 = 60$, günstig 6: $P = \frac{6}{60} = \frac{1}{10} = 10\,\%$ (nicht $\frac{6}{54}$)}
\lsg{3.}{$44\,\%$}{alle $11 + 6 + 8 = 25$, günstig 11: $P = \frac{11}{25} = 0{,}44 = 44\,\%$}
\lsg{4.}{$\frac{1}{8} = 12{,}5\,\%$}{alle $600 - 201 + 1 = 400$ Lose, günstig $250 - 201 + 1 = 50$: $P = \frac{50}{400} = \frac{1}{8} = 0{,}125 = 12{,}5\,\%$}
\lsg{5.}{Bude A: $25\,\%$ gegen $22\,\%$}{Bude A. A: alle $150 + 50 = 200$, $P = \frac{50}{200} = 25\,\%$; B: $P = \frac{66}{300} = 22\,\%$. Falsch wäre $\frac{50}{150}$.}
\par\Needspace{25mm}\vspace{6pt}\noindent\colorbox{black!12}{\makebox[7mm]{\bfseries\strut C}}\hspace{2mm}{\bfseries Gegenereignis: „nicht“ und „weder … noch“}\par\vspace{4pt}
\lsg{1.}{$0{,}7$ (also $70\,\%$)}{$1 - 0{,}3 = 0{,}7$}
\lsg{2.}{$P = \frac{5}{6}$}{$P(6) = \frac{1}{6}$; $P(\text{keine Sechs}) = 1 - \frac{1}{6} = \frac{5}{6}$}
\lsg{3.}{$P = \frac{2}{3}$}{alle 12; $P(\text{grün}) = \frac{4}{12}$; $P(\text{nicht grün}) = 1 - \frac{4}{12} = \frac{8}{12} = \frac{2}{3}$}
\lsg{4.}{$P = \frac{1}{2}$}{durch 3 oder durch 4 teilbar: 3, 4, 6, 8, 9, 12 – das sind 6 (die 12 nur einmal zählen); $P = 1 - \frac{6}{12} = \frac{6}{12} = \frac{1}{2}$}
\lsg{5.}{Nein: Ole $\frac{2}{3}$, Mia nur $\frac{1}{3}$}{Nein. $P(\text{Mia}) = P(1 \text{ oder } 2) = \frac{2}{6} = \frac{1}{3}$; $P(\text{Ole}) = 1 - \frac{1}{3} = \frac{2}{3}$. Ole gewinnt doppelt so oft wie Mia.}
\par\Needspace{25mm}\vspace{6pt}\noindent\colorbox{black!12}{\makebox[7mm]{\bfseries\strut D}}\hspace{2mm}{\bfseries Wenn sich die Grundmenge ändert}\par\vspace{4pt}
\lsg{1.}{$P = \frac{1}{3}$}{noch $10 - 1 = 9$ Kugeln, rot 3: $P = \frac{3}{9} = \frac{1}{3}$}
\lsg{2.}{$P = \frac{4}{11}$}{noch $12 - 1 = 11$ Kugeln, davon $5 - 1 = 4$ gelb: $P = \frac{4}{11}$}
\lsg{3.}{$P = \frac{3}{25} = 12\,\%$}{möglich sind nur noch die Nummern mit Endziffer 3: 103, 113, \ldots, 343 – das sind 25. Davon enden auf 13: 113, 213, 313, also 3. $P = \frac{3}{25} = 0{,}12 = 12\,\%$ (nicht $\frac{3}{250}$)}
\lsg{4.}{Nein: $\frac{4}{50} = 8\,\%$}{Nein. Lose mit Endziffer 7: 307, 317, \ldots, 797 – das sind 50. Davon enden auf 27: 327, 427, 527, 627, 727; ohne 527 bleiben 4. $P = \frac{4}{50} = 0{,}08 = 8\,\%$, also weniger als 10\,\%.}
\par\Needspace{25mm}\vspace{6pt}\noindent\colorbox{black!12}{\makebox[7mm]{\bfseries\strut E}}\hspace{2mm}{\bfseries Zufallsgerät bauen und vorhersagen}\par\vspace{4pt}
\lsg{1.}{2 rote Felder}{$\frac{1}{4} \cdot 8 = 2$ Felder}
\lsg{2.}{8 weiße Kugeln}{Die 4 schwarzen sind $\frac{1}{3}$ aller Kugeln, also alle $3 \cdot 4 = 12$ Kugeln; weiße: $12 - 4 = 8$}
\lsg{3.}{$\approx 0{,}173$ gegen $\approx 0{,}167$: passt gut}{relative Häufigkeit $\frac{52}{300} \approx 0{,}173$; $P(6) = \frac{1}{6} \approx 0{,}167$. Fast gleich: Bei vielen Würfen liegt die relative Häufigkeit nahe an der Wahrscheinlichkeit.}
\lsg{4.}{Nein: etwa 60 nötig, rund 10 zu wenig}{Nein. $P = \frac{2}{16} = \frac{1}{8}$; erwartet $\frac{1}{8} \cdot 480 = 60$ Teddys; es fehlen etwa $60 - 50 = 10$.}
\par\Needspace{25mm}\vspace{6pt}\noindent\colorbox{black!12}{\makebox[7mm]{\bfseries\strut T}}\hspace{2mm}{\bfseries Probetest}\par\vspace{4pt}
\lsg{1.}{$\frac{1}{2}$}{4 von 8 Feldern sind rot, $\frac{4}{8} = \frac{1}{2}$. Nicht $\frac{1}{3}$ – es zählen die Felder, nicht die Farben.}
\lsg{2.}{$\frac{2}{15} \approx 13{,}3\,\%$}{alle $104 + 16 = 120$, günstig 16: $P = \frac{16}{120} = \frac{2}{15} \approx 13{,}3\,\%$}
\lsg{3.}{$P = \frac{3}{4}$}{über 9: 10, 11, 12 – das sind 3; $P = 1 - \frac{3}{12} = \frac{9}{12} = \frac{3}{4}$}
\lsg{4.}{Nein: $\frac{3}{12} = \frac{1}{4}$}{Nein. Es sind nur noch $15 - 3 = 12$ Bonbons, davon 3 mit Zitrone: $P = \frac{3}{12} = \frac{1}{4}$.}
\lsg{5.}{3 Felder; etwa 75 nötig – 70 reichen nicht}{$0{,}25 \cdot 12 = 3$ Preisfelder; erwartet $0{,}25 \cdot 300 = 75$ Preise. Nein, 70 reichen wahrscheinlich nicht; es fehlen etwa $75 - 70 = 5$.}
\end{document}
